\documentclass[aspectratio=169,11pt]{beamer}

\usepackage[T1]{fontenc}
\usepackage[utf8]{inputenc}
\usepackage{lmodern}
\usepackage{amsmath,amssymb}
\usepackage{booktabs}
\usepackage{xcolor}
\usepackage{hyperref}
\usepackage{microtype}

\usetheme{Madrid}
\usecolortheme{default}

\definecolor{LaurierPurple}{HTML}{4B116F}
\definecolor{DeepBlue}{HTML}{0047AB}

\setbeamertemplate{navigation symbols}{}
\setbeamercolor{title}{fg=white,bg=LaurierPurple}
\setbeamercolor{frametitle}{fg=white,bg=LaurierPurple}
\setbeamercolor{block title}{fg=white,bg=DeepBlue}
\setbeamercolor{structure}{fg=DeepBlue}

\hypersetup{colorlinks=true,urlcolor=DeepBlue,linkcolor=DeepBlue}

\title[Supervised \& Unsupervised]{Module 2: Supervised and Unsupervised Learning in Economics}
\subtitle{EC 410K / EC 610I --- Machine Learning in Economics}
\author{M. Jahangir Alam, PhD}
\institute{Wilfrid Laurier University}
\date{Spring 2026}

\begin{document}

\begin{frame}[plain]
  \titlepage
\end{frame}

\begin{frame}{Module overview}
  \begin{itemize}
    \item \textbf{Supervised learning:} learn \(\widehat f\) mapping features \(X\) to outcome \(Y\) using labeled data.
    \item \textbf{Unsupervised learning:} discover structure in \(X\) (clusters, low-dimensional representations).
    \item Economics adds: \textbf{interpretation}, \textbf{stability}, and separating \textbf{prediction} from \textbf{causal} claims.
  \end{itemize}
\end{frame}

\begin{frame}{Learning objectives}
  \begin{enumerate}
    \item Explain train/validation/test evaluation and why in-sample fit misleads.
    \item Implement regularized regression and a flexible learner (random forest) with cross-validation.
    \item Use PCA and k-means and discuss what they do (and do not) imply economically.
    \item Connect methods to Mullainathan and Spiess (2017) and Gu, Kelly, and Xiu (2020).
  \end{enumerate}
\end{frame}

\begin{frame}{Background: why economists care}
  \begin{itemize}
    \item Many problems are naturally \textbf{high-dimensional}: many predictors, interactions, latent types.
    \item Classic linear models may miss nonlinearities; ML methods may improve \textbf{predictive accuracy}.
    \item But improved prediction does not automatically answer causal policy questions (Module 1).
  \end{itemize}
\end{frame}

\begin{frame}{Supervised learning setup}
  Given pairs \((Y_i, X_i)\), choose \(\widehat f\) in a class \(\mathcal{F}\):
  \[
    \widehat f \in \arg\min_{f \in \mathcal{F}} \;\frac{1}{n}\sum_{i=1}^n \ell\big(Y_i, f(X_i)\big) + \text{complexity penalty}
  \]
  \begin{itemize}
    \item Regression losses: squared error, Huber; classification: log loss, hinge; etc.
  \end{itemize}
\end{frame}

\begin{frame}{Overfitting and the bias--variance tradeoff (intuition)}
  \begin{itemize}
    \item \textbf{Bias:} systematic error from wrong functional form / restrictions.
    \item \textbf{Variance:} sensitivity to training noise; tends to rise with model flexibility.
    \item Goal: good \textbf{out-of-sample} performance, not maximizing in-sample \(R^2\).
  \end{itemize}
\end{frame}

\begin{frame}{Cross-validation (workflow)}
  \begin{enumerate}
    \item Split data into \(K\) folds; rotate which fold is held out.
    \item Tune hyperparameters using average holdout loss.
    \item Report performance honestly (avoid ``testing on the training set'').
  \end{enumerate}
  \vspace{0.3em}
  \begin{alertblock}{Economics note}
    Time series need specialized splitting (rolling windows), not i.i.d.\ folds.
  \end{alertblock}
\end{frame}

\begin{frame}{Gu, Kelly, and Xiu (2020): empirical asset pricing via ML}
  \begin{itemize}
    \item Massive predictor sets in asset pricing: compare many ML methods for \textbf{return prediction}.
    \item Central lesson for this course: careful out-of-sample design and economic interpretation of signals.
    \item Students are not required to replicate the full paper; the Colab exercise captures the \textbf{spirit}: many \(Z\), compare learners.
  \end{itemize}
\end{frame}

\begin{frame}{Mullainathan and Spiess (2017): ML as applied econometrics}
  \begin{itemize}
    \item Bridges ML practice with econometric goals: prediction, estimation, hypothesis formation.
    \item Emphasizes discipline about goals (prediction vs structure) and validation.
  \end{itemize}
\end{frame}

\begin{frame}{Unsupervised learning: PCA (idea)}
  \begin{itemize}
    \item Find directions of maximum variance in standardized \(X\): scores \(Z_k^\top X\).
    \item Useful for visualization, compression, and building features---but PCs are not automatically ``economic factors.''
  \end{itemize}
\end{frame}

\begin{frame}{Unsupervised learning: k-means (idea)}
  \begin{itemize}
    \item Partition units into \(K\) groups minimizing within-cluster dispersion.
    \item Clusters can be useful stylized facts or segmentation---but labeling requires domain reasoning.
  \end{itemize}
\end{frame}

\begin{frame}{Connection to applied research: ADH (2013) as a discussion example}
  \begin{itemize}
    \item \textbf{Supervised angle:} predict local manufacturing decline from pre-shock covariates (``vulnerability'').
    \item \textbf{Unsupervised angle:} cluster commuting zones by industry mix / shock exposure geometry.
    \item \textbf{Causal angle:} separate from prediction---requires the paper’s identification logic.
  \end{itemize}
\end{frame}

\begin{frame}{Python / Colab demo plan}
  \begin{enumerate}
    \item Simulate many characteristics predicting returns (sparse linear truth + noise).
    \item Compare Ridge vs random forest using 5-fold CV MSE.
    \item Run PCA + k-means; plot PC1--PC2; describe clusters descriptively.
  \end{enumerate}
\end{frame}

\begin{frame}[fragile]{Colab setup}
  \begin{verbatim}
!pip -q install numpy pandas scikit-learn matplotlib
  \end{verbatim}
\end{frame}

\begin{frame}{Student / group assignment}
  \begin{itemize}
    \item Present MS (2017) + GKX (2020): question, data setting (GKX), main ML takeaway.
    \item Teach cross-validation clearly with a diagram on the board or slides.
    \item Run \texttt{colab.py}; explain each output block in economics language.
    \item Propose one extension idea for ADH or another replication paper (prediction feature, clustering, risk score).
  \end{itemize}
\end{frame}

\begin{frame}{Summary}
  \begin{itemize}
    \item Supervised ML optimizes predictive performance with disciplined out-of-sample checks.
    \item Unsupervised methods summarize \(X\); economic meaning is not automatic.
    \item Connect every algorithmic output to a \textbf{question} and a \textbf{validation strategy}.
  \end{itemize}
\end{frame}

\begin{frame}[plain]
  \begin{center}
    \Large Questions?\\[0.8em]
    \normalsize Next: linear models and regularization (Module 3).
  \end{center}
\end{frame}

\end{document}
